%==============================================================================%
% Chapter 11 – End-to-End Correctness
%==============================================================================%

\chapter{End-to-End Correctness}
\label{ch:end_to_end}

\section{Overview of \texttt{NTTEndToEnd.ec}}

The file \texttt{NTTEndToEnd.ec} (144 lines) is the apex of the verification
chain.  It imports all preceding modules and states the final correctness
theorems that directly relate the exported \jasmin{} functions to the abstract
spectral NTT specification.

\begin{lstlisting}[language=EasyCrypt,
  caption={Module imports in \texttt{NTTEndToEnd.ec}},
  label={lst:e2e_imports}]
require import NTTFullSpec NTTFullAlgebra.
require import NTT_Fq Fq GFq Rq.
require import RefJasminNTTLoop CTLoopEquiv.
require import Hpoly_extract Hpoly_loop.
\end{lstlisting}

\section{Main Correctness Theorems}

\subsection{Forward NTT Correctness}

\begin{theorem}[Jasmin NTT correctness]
  \label{thm:main_ntt}
  Let \(\texttt{rp}\) be a Jasmin word array represented in EasyCrypt as a
  \ec{BArray1024.t}, encoding a polynomial \(p\) with 16-bit signed
  coefficients.  Then executing \jasm{poly\_ntt\_jazz(\texttt{rp})} produces
  an output array $\texttt{rp}'$ satisfying:
  \[
    \texttt{rp}' \;\text{represents}\; \op{full\_ntt}(p)
    \;\text{with 24-bit bounds},
  \]
\end{theorem}

\begin{lstlisting}[language=EasyCrypt,
  caption={End-to-end NTT correctness statement},
  label={lst:e2e_ntt}]
op jasmin_ntt_full_post (rp : BArray1024.t) (p : poly) =
  NTT_Fq.poly_repr_bound rp (NTTFullSpec.full_ntt p) 24.

lemma haetae_jasmin_ntt_correct p :
  hoare [Hpoly_extract.M.poly_ntt_jazz :
    NTT_Fq.poly_repr_bound rp p 16 ==>
    NTT_Fq.poly_repr_bound res (NTTFullSpec.full_ntt p) 24].
\end{lstlisting}

\subsection{Inverse NTT Correctness}

\begin{theorem}[Jasmin inverse NTT correctness]
  \label{thm:main_invntt}
  Let $\texttt{rp}$ encode a polynomial \(p\) satisfying the 16-bit bound
  required by the checked inverse kernel theorem.
  Then executing \jasm{poly\_invntt\_jazz(\texttt{rp})} produces an output
  array $\texttt{rp}'$ satisfying:
  \[
    \texttt{rp}' \;\text{represents}\;
      \ec{array256\_mont}(\op{full\_invntt}(p)),
  \]
  i.e. the inverse output is represented with one Montgomery factor \(R\).
\end{theorem}

\begin{lstlisting}[language=EasyCrypt,
  caption={End-to-end inverse NTT correctness statement},
  label={lst:e2e_invntt}]
op jasmin_invntt_full_post (rp : BArray1024.t) (p : poly) =
  NTT_Fq.poly_repr_bound rp
    (NTT_Fq.array256_mont (NTTFullSpec.full_invntt p)) 16.

lemma haetae_jasmin_invntt_correct p :
  hoare [Hpoly_extract.M.poly_invntt_jazz :
    NTT_Fq.poly_repr_bound rp p 16 ==>
    NTT_Fq.poly_repr_bound res
      (NTT_Fq.array256_mont (NTTFullSpec.full_invntt p)) 16].
\end{lstlisting}

The \ec{array256\_mont} postcondition is the precise Montgomery-form
convention: each abstract coefficient is multiplied by \(R\), not by
\(R^{-1}\).  The value \(R^{-1}\) appears internally when cancelling a
Montgomery multiplication, but it is not the final representation factor.

\section{Round-Trip Consequence}

\begin{proposition}[Abstract round-trip with Montgomery output]
  \label{thm:roundtrip}
  For any polynomial $a \in \Z_q^{256}$:
  \[
    \op{full\_invntt}(\op{full\_ntt}(a)) = a.
  \]
Consequently, a concrete inverse call whose input represents
\(\op{full\_ntt}(a)\) and satisfies the inverse theorem's representation
bound returns a word array representing \(\ec{array256\_mont}(a)\), i.e.,
\(a\cdot R\) componentwise.
\end{proposition}

\begin{proof}
  At the mathematical layer:
  \begin{align*}
    \op{full\_invntt}(\op{full\_ntt}(a))[i]
    &= a[i],
  \end{align*}
  by \Cref{thm:intt_correct}.  The concrete inverse theorem then wraps the
  result in \ec{array256\_mont}, giving the \(R\)-scaled representation.
\end{proof}

\Cref{thm:roundtrip} is a representation-level consequence, not an additional
top-level theorem in \ec{NTTEndToEnd.ec}.  The checked file states the forward
and inverse exported-function theorems separately, and the dissertation does
not cite a direct one-procedure in-place pipeline theorem.

\section{Proof Composition}

The end-to-end proofs are assembled from four chains:

\begin{enumerate}
  \item \textbf{Full-safe extraction $\to$ proof loop}:
        \texttt{CTLoopEquiv.ec} proves that the helperized full-safe
        \texttt{Hpoly\_extract} procedures implement the same transition as
        the proof-only direct-loop model in \texttt{Hpoly\_loop.ec}.
  \item \textbf{Proof loop $\to$ Imperative}: \Cref{thm:ntt_loop_equiv} is
        instantiated by \texttt{RefJasminNTT\allowbreak Loop.ec}.
  \item \textbf{Imperative $\to$ Abstract}: \Cref{thm:ntt_correct} from
        \texttt{NTTFullAlgebra.ec}.
  \item \textbf{Abstract inverse identity}: \Cref{thm:intt_correct} is the
        pen-and-paper field argument stated in this dissertation.  The
        checked EasyCrypt files prove the forward and inverse maps against
        \ec{full\_ntt} and \ec{full\_invntt} separately; they do not add a
        standalone checked round-trip lemma.
\end{enumerate}

The composition is expressed in \easycrypt{} as sequential lemma application:

\begin{lstlisting}[language=EasyCrypt,
  caption={Proof composition in \texttt{NTTEndToEnd.ec}},
  label={lst:e2e_proof}]
proof haetae_jasmin_ntt_correct.
  move => rp a ha.
  (* Step 1: proof-only core loop -> imperative -> abstract. *)
  have hloop_core :
    hoare [Hpoly_loop.M._poly_ntt :
      NTT_Fq.poly_repr_bound rp a 16 ==>
      jasmin_ntt_full_post res a].
  + by conseq RefJasminNTT.poly_ntt_core_ref
      (NTTFullAlgebra.ntt_full_ntt a) => /#.
  (* Step 2: proof-only exported wrapper. *)
  have hloop_jazz :
    hoare [Hpoly_loop.M.poly_ntt_jazz :
      NTT_Fq.poly_repr_bound rp a 16 ==>
      jasmin_ntt_full_post res a].
  + by conseq RefJasminNTT.poly_ntt_jazz_ref hloop_core => /#.
  (* Step 3: helperized full-safe extraction -> proof-only wrapper. *)
  by conseq CTLoopEquiv.poly_ntt_jazz_ct_loop hloop_jazz => /#.
qed.
\end{lstlisting}

\section{Status of the Proof}

\subsection{Completed Machine-Checked Steps}

\begin{enumerate}
  \item \textbf{\texttt{Fq.ec}}: All lemmas proved, including Montgomery
        reduction correctness, signed arithmetic identities, and bit-width
        bound propagation.

  \item \textbf{\texttt{NTT\_Fq.ec}}: Full imperative specification including
        termination proofs and loop invariants.

  \item \textbf{\texttt{NTTFullSpec.ec}}: Abstract NTT definitions and
        conversion lemmas from the relational specifications to
        \ec{full\_ntt} and \ec{full\_invntt}.

  \item \textbf{\texttt{RefJasminNTT.ec}}: The preserved original direct-loop
        proof boundary and its two main refinement lemmas
        \ec{poly\_ntt\_core\_ref} and \ec{poly\_invntt\_core\_ref}.

  \item \textbf{\texttt{Hpoly\_loop.ec}}: The proof-only direct-loop model
        used as the stable target for the full-safe helperized extraction.

  \item \textbf{\texttt{RefJasminNTT\allowbreak Loop.ec}}: The direct-loop refinement
        proof retargeted to \texttt{Hpoly\_loop}.

  \item \textbf{\texttt{CTLoopEquiv.ec}}: The pRHL bridge proving equivalence
        between the full-safe \texttt{Hpoly\_extract} procedures and
        \texttt{Hpoly\_loop}, including exported wrappers.

  \item \textbf{\texttt{NTTFullAlgebra.ec}}: The stage-induction theorems
        \ec{ntt\_full\_ntt} and \ec{invntt\_full\_invntt}.

  \item \textbf{\texttt{NTTEndToEnd.ec}}: The composed top-level Hoare
        statements for the exported Jasmin functions.
\end{enumerate}

\subsection{Machine-Checked Top-Level Compilation}

The checked development currently compiles through \texttt{NTTEndToEnd.ec} in
\texttt{easycrypt-ct}.  This means the full-safe helperized Jasmin extraction,
the proof-loop refinement, the imperative-to-abstract stage proofs, and the
final composition are all discharged for the scalar NTT and inverse NTT
covered by this dissertation.  Remaining limitations are scope limitations,
not open proof gaps in the NTT chain: AVX2 vectorisation, base multiplication,
and full protocol-level security integration are outside the proved artifact.

\section{Summary of Verification Metrics}

\begin{table}[htbp]
\centering
\caption{Verification metrics by file}
\label{tab:metrics}
\small
\begin{tabular}{
  >{\raggedright\arraybackslash}p{0.32\textwidth}
  rr
  >{\raggedright\arraybackslash}p{0.28\textwidth}}
\toprule
File & Lines & Named lemmas & Status \\
\midrule
\texttt{Fq.ec}           & 535  & 16  & Compiled \\
\texttt{GFq.ec}          & 27   & 0   & Compiled field instantiation \\
\texttt{Rq.ec}           & 117  & 0   & Compiled ring interface \\
\texttt{Montgomery.ec}   & 427  & 29  & Compiled \\
\texttt{NTT\_Fq.ec}      & 988  & 52  & Compiled \\
\texttt{NTTFullSpec.ec}  & 160  & 10  & Compiled \\
\texttt{NTTFullAlgebra.ec} & 3085 & 111 & Compiled \\
\texttt{RefJasminNTT.ec} & 2609 & 103 & Preserved direct-loop proof \\
\texttt{Hpoly\_loop.ec} & 156 & 0 & Compiled proof-loop model \\
\texttt{RefJasminNTT\allowbreak Loop.ec} & 2629 & 105 & Compiled full-safe refinement target \\
\texttt{CTLoopEquiv.ec} & 657 & 14 & Compiled helper bridge \\
\texttt{NTTEndToEnd.ec}  & 144  & 10  & Compiled full-safe top-level \\
\midrule
\textbf{Total}           & \textbf{11534} & \textbf{450} & \textbf{Compiled through full-safe end-to-end} \\
\bottomrule
\end{tabular}
\normalsize
\end{table}
